\chapter{Enlargeability}
\label{chapter:enlargeability}

In this chapter we discuss the concept of enlargeability, which is the central idea that the results in this thesis are concerned with. The following definition is, for example, adopted in \cite{LawsonMichelsohn1989SpinGeometry, BrunnbauerHanke, cecchini2021enlargeable}.
\begin{definition}
    \label{enl:def:enlargeability}
    A closed oriented Riemannian $n$-manifold $M$ is called enlargeable if for all $\epsilon >0$ there exists a Riemannian covering $\hat{M} \to M$ together with an $\epsilon$-Lipschitz map $\hat{M} \to S^n$, which is constant outside a compact subset and of non-zero degree.
\end{definition}

Note that this notion does not depend on the particular choice of Riemannian metric on $M$, since for closed manifolds all metrics are in bi-Lipschitz correspondence. Intuitively, the covering $\hat{M}$ can be pictured as an ``expansion'' of $M$ itself (think, for example, of the usual sketches of coverings of $S^1$). Enlargeability now means that this expansion can be chosen arbitrarily large, so that it can be wrapped as tightly around a sphere as one wishes. 
The simplest example of an enlargeable manifold is the $n$-torus $T^n$. To see this consider the cover $\R^n \to T^n$. Clearly $\R^n$ has the property that for any $\epsilon > 0$ there exists an $\epsilon$-Lipschitz map $\R^n \to S^n$, which constant outside a compact subset and of non-zero degree. We refer to manifolds with this property as \emph{hyperspherical}. More generally, if a closed oriented manifold $M$ admits a metric of non-positive sectional curvature, its universal cover is $\R^n$ by the Cartan-Hadamard Theorem and consequently $M$ is enlargeable. Non-enlargeable manifolds are, for example, all simply connected manifolds, since all covers are trivial in this case.
In view of the above examples, one might naively expect that the universal cover of an enlargeable manifold - which in some sense is the ``largest'' covering - is always hyperspherical. However, Brunnbauer and Hanke have constructed counterexamples to this naive expectation \cite[Thm. 1.4]{BrunnbauerHanke}.

The central fact about enlargeable manifolds is the following theorem.

\begin{theorem}
    \label{enl:thm:enl_mnf_do_not_admit_psc}
    An enlargeable manifold does not admit a metric of positive scalar curvature.
\end{theorem}

This is the result that originally motivated the definition of enlargeability by Gromov and Lawson \cite{gromov1980spin}, where they first showed that enlargeability is an obstruction to the existence of a \psc-metric, under the additional assumption that the coverings $\hat{M}$ in Definition \ref{enl:def:enlargeability} are finite and spin. Note that $T^n$ is still enlargeable even under these additional assumptions. In fact, using their result about enlargeability, Gromov and Lawson gave the first general proof of the fact that $T^n$ does not admit a \psc-metric for any dimension $n$. In a later paper, they were able to drop the finiteness assumption on $\hat{M}$ \cite{Gromov1983}. Both proofs are based on techniques involving the Dirac operator, which requires the existence of a spin structure, as briefly discussed in \hyperref[chapter:psc]{Chapter 2}. 
So it was unclear whether and how these results carry over to the non-spin case. Rather recently, however, Cecchini and Schick resolved this issue \cite{cecchini2021enlargeable}. Their proof is based on another fundamental method for obtaining obstructions to positive scalar curvature: \emph{minimal hypersurfaces}. This approach is due to Schoen and Yau, and the core observation is the following. Suppose a closed manifold $M^n$ carries a \psc-metric and $N^{n-1} \subseteq M$ is a two-sided hypersurface, which is a local minimum of the volume functional. Then $N$ also admits a \psc-metric \cite[Proof of Thm. 1]{SchoenYau1979}, \cite[Prop. 2.25]{Lee2019}. So, to obstruct the existence of a \psc-metric on $M$, we can try to obstruct the existence of a \psc-metric on $N$. This idea can then be used inductively to reduce the problem to one of low dimension, where more methods are available (for example, the Gauß-Bonnet Theorem in dimension $2$). The problem with this approach is the existence of such a minimal hypersurface $N$. In general, existence results for minimal hypersurfaces only hold up to dimension $n \leq 7$. However, at least in certain cases, Schoen and Yau proposed methods for overcoming this dimensional restriction \cite{schoen2017positivescalarcurvatureminimal}, which is what Cecchini and Schick ultimately relied on to extend Theorem \ref{enl:thm:enl_mnf_do_not_admit_psc} to the non-spin case.

The following lemma now provides two simple methods to construct new enlargeable manifolds out of given ones. 

\begin{lemma}
    \label{enl:lem:constructing_enlargeable_mnf}
    \begin{enumerate}[label=(\roman*)]
        \item The product of enlargeable manifolds is enlargeable. \label{enl:eq:lem_for_constructing_enlargeable_mnf_1}
        \item A closed oriented manifold admitting a non-zero degree map onto an enlargeable manifold is enlargeable. \label{enl:eq:lem_for_constructing_enlargeable_mnf_2}
    \end{enumerate}
\end{lemma}

\begin{proof}
    For \ref{enl:eq:lem_for_constructing_enlargeable_mnf_1}, let $M^m$ and $N^n$ be enlargeable. For $\epsilon > 0$, choose Riemannian coverings $p \colon \hat{M} \to M$ and $q \colon \hat{N} \to N$ as well as $\epsilon$-Lipschitz maps $f \colon \hat{M} \to S^m$ and $g \colon \hat{N} \to S^n$, which are constant outside a compact subset and of non-zero degree. Then 
    \[
    p \times q \colon \hat{M} \times \hat{N} \to M \times N
    \]
    is a Riemannian covering, where both manifolds are equipped with the corresponding product metric and the map
    \[
    f \times g \colon \hat{M} \times \hat{N} \to S^m \times S^n
    \]
    is $\epsilon$-Lipschitz. Let $y \in S^m$ and $z \in S^n$ denote the respective constant values of $f$ and $g$ outside compact subsets. Consider the map 
    \[
    \psi \colon S^m \times S^n \to S^m \wedge S^n \cong S^{m+n},
    \]
    given by projection onto the smash product 
    \[
    S^m \wedge S^n := \raisebox{3pt}{$\modulo{S^m \times S^n}{(S^m \times \{z\} \cup \{y\} \times S^n)}$},
    \]
    which is well known to be homeomorphic to the sphere $S^{m+n}$. Then the composition 
    \[
    \psi \circ (f \times g) \colon \hat{M} \times \hat{N} \to S^{m+n}
    \] 
    is again constant outside a compact subset and has degree $\mathrm{deg}(f) \mathrm{deg}(g) \neq 0$. After a slight perturbation me may assume $\psi$ to be smooth, so that it is $\lambda$-Lipschitz for some $\lambda > 0$, since $S^m \times S^n$ is compact. Consequently $\psi \circ (f \times g)$ is $\lambda \epsilon$-Lipschitz, and since $\epsilon$ was chosen arbitrarily, the statement follows.

    For \ref{enl:eq:lem_for_constructing_enlargeable_mnf_2}, let $M^m$ be enlargeable and let $h \colon P \to M$ be some non-zero degree map for a closed oriented manifold $P$. Again, given $\epsilon > 0$, choose a Riemannian cover $p \colon \hat{M} \to M$ as well as an $\epsilon$-Lipschitz map $f \colon \hat{M} \to S^m$. Define a cover $\pi \colon \hat{P} \to P$ as the pullback cover of $p$ under $h$, that is, $\pi$ fits into a pullback diagram of the form 
    \[\begin{tikzcd}
	{\hat{P}} & {\hat{M}} \\
	P & M
	\arrow["H", from=1-1, to=1-2]
	\arrow["\pi", from=1-1, to=2-1]
	\arrow["p", from=1-2, to=2-2]
	\arrow["h", from=2-1, to=2-2]
    \end{tikzcd}\]
    for some smooth map $H$. Fix metrics on $P$ and $\hat{P}$, such that $\pi$ is a Riemannian cover. Then $h$ is $\mu$-Lipschitz for some constant $\mu > 0$, since $P$ is closed. Consequently, $H$ is also $\mu$-Lipschitz, given that both $\pi$ and $p$ are local isometries. Now consider the map 
    \[
    f \circ H \colon \hat{P} \to S^m.
    \] 
    From the above discussion it follows that this map is $\mu \epsilon$-Lipschitz. Furthermore, since $h$ is proper and $H$ is an isomorphism on fibers, we conclude that $H$ is proper as well. Therefore $f \circ H$ is constant outside a compact subset. Counting preimages of a regular value one sees that 
    \[
    \mathrm{deg}(f \circ H) = \mathrm{deg}(h) \mathrm{deg}(f) \neq 0.
    \]
\end{proof}

From Part \ref{enl:eq:lem_for_constructing_enlargeable_mnf_2} of the above lemma it follows immediately that enlargeability is a homotopy invariant, since a homotopy equivalence of connected closed oriented manifolds has degree $\pm 1$. In fact, much more is true: enlargeability is a homological invariant in the following sense. 

\begin{theorem}[{\cite[Thm. 1.3]{BrunnbauerHanke}}]
    Let $M$ be a closed oriented $n$-manifold and let $\phi \colon M \to B\pi_1(M)$ be the classifying map of the universal cover. Then there exists a rational vector subspace 
    \[
    H_n^\mathrm{small}(B\pi_1(M);\Q) \subseteq H_n(B\pi_1(M);\Q)
    \]
    of ``small'' classes with the following property. $M$ is enlargeable if and only if the image $\phi_\ast [M]$ of the rational fundamental class $[M] \in H_n(M;\Q)$ of $M$ is not contained in  $H_n^\mathrm{small}(B\pi_1(M);\Q)$. In other words $\phi_\ast [M]$ is a ``large'' class.
\end{theorem}

 In \cite{BrunnbauerHanke}, the authors gave a general definition of ``small'' and ``large'' homology classes of a simplicial complex $X$ with finitely generated fundamental group. With this definition the assignment $X \mapsto H_n^\mathrm{small}(X;\Q)$ is functorial for continuous maps $f \colon X \to Y$ inducing an epimorphism on $\pi_1$, that is, $f_\ast \colon H_n(X;\Q) \to H_n(Y;\Q)$ maps small classes to small classes. Moreover, if $f$ happens to induce an isomorphism on $\pi_1$, then $f_\ast$ also maps large classes to large classes. This general definition then specializes to the fact that a closed oriented manifold $M$ is enlargeable if and only if the fundamental class $[M]$ is large, that is, $[M] \notin H_n^\mathrm{small}(M;\Q)$. Combining these observations, the theorem follows, since the classifiying map $\phi \colon M \to B\pi_1(M)$ always induces an isomorphism on $\pi_1$. In practice one can for example use this theorem to obstruct enlargeability of $M$ by showing that $\phi_\ast [M] = 0$. More importantly, the argument outlined above establishes that the notion of enlargeability can be understood purely in terms of simplicial complexes. 
